\section{正规方程解多元线性回归}
\label{sec:chapter3}

\subsection{正规方程}

正规方程利用多元微分直接求解线性回归的最优参数，无须像梯度下降那样反复迭代。

对于代价函数：
\[
J(\theta)=J(\theta_0, \theta_1,\cdots, \theta_n)
\]
若代价函数可微且为凸函数，则可令各偏导数为零，从而求得使其取最小值的参数 $\theta$：
\[
\frac{\partial J}{\partial \theta_j}=0,\quad j=0,1,\cdots,n
\]
写成向量形式即为：
\[
\frac{\partial J}{\partial \theta}=\vec 0
\]
下面将这一方法应用于多元线性回归。

多元线性回归的代价函数为：
\[
J(\theta)=\frac{1}{2m}\sum_{i=1}^m\left(\theta^Tx^{(i)}-y^{(i)}\right)^2
\]
其梯度为：
\[
\frac{\partial J}{\partial \theta}=\frac{1}{m}\sum_{i=1}^m\left(\theta^Tx^{(i)}-y^{(i)}\right)x^{(i)}
\]
令其为零，解得：
\[
\theta^Tx^{(i)}=y^{(i)},\quad i=1,2,\cdots,m
\]
这些等式可以统一写成矩阵形式：
\[
X\theta=y
\]
其中，
\[
X=\begin{bmatrix}x_0^{(1)}&x_1^{(1)}&\cdots& x_n^{(1)}\\x_0^{(2)}&x_1^{(2)}&\cdots& x_n^{(2)}\\\vdots&\vdots&\ddots&\vdots\\x_0^{(m)}&x_1^{(m)}&\cdots& x_n^{(m)}\end{bmatrix}=\begin{bmatrix}{x^{(1)}}^T\\{x^{(2)}}^T\\\vdots\\{x^{(m)}}^T\\\end{bmatrix},\quad y=\begin{bmatrix}y^{(1)}\\y^{(2)}\\\vdots\\y^{(m)}\end{bmatrix}
\]
两边同时乘以 $X^T$，并假设 $X^TX$ 可逆，解得：
\[
\theta=(X^TX)^{-1}X^Ty
\]
这就是多元线性回归的解析解。

由于 $X^TX$ 是一个 $(n+1)\times(n+1)$ 的矩阵，直接求逆或求解线性方程组的时间复杂度通常为 $O(n^3)$. 当 $n$ 不大时，正规方程能够一次得到解析解，也不需要选择学习率；但当 $n$ 达到 $10^4,10^5$ 或更高数量级时，矩阵运算的时间和空间开销都会迅速增加，此时通常应考虑梯度下降等迭代方法。


\subsection{代码实现}

仍使用第 \ref{sec:chapter2} 节中的多元线性回归数据，代码如下：

\begin{lstlisting}[language=Python]
import numpy as np

def J(T, X, Y):
	res = 0
	for i in range(m):
		res += (np.matmul(T.T, X[i:i+1, :].T) - Y[i:i+1, :]) ** 2
	res /= 2 * m;
	return res

data = np.genfromtxt("ex1data2.txt", delimiter = ',')
(m, n) = data.shape
X = np.column_stack((np.ones((m, 1)), data[:, :-1]))
Y = data[:, -1:]
T = np.matmul(np.matmul(np.linalg.inv(np.matmul(X.T, X)), X.T), Y)
print(T)
print(J(T, X, Y))
\end{lstlisting}

计算很快即可完成，结果为：$\theta=(89597.9095428,\,139.21067402,\,-8738.01911233)^T$.


\subsection{不可逆情形}

上述求解过程建立在 $X^TX$ 可逆的假设上。当特征之间存在冗余或线性相关关系时，该矩阵可能不可逆。此时可以将代码中的 \texttt{inv()} 替换为 \texttt{pinv()}，通过\textbf{伪逆矩阵}求得一个最小二乘意义下的解。

\begin{lstlisting}[language=Python]
import numpy as np

def J(T, X, Y):
	res = 0
	for i in range(m):
		res += (np.matmul(T.T, X[i:i+1, :].T) - Y[i:i+1, :]) ** 2
	res /= 2 * m;
	return res

data = np.genfromtxt("ex1data2.txt", delimiter = ',')
(m, n) = data.shape
X = np.column_stack((np.ones((m, 1)), data[:, :-1]))
Y = data[:, -1:]
T = np.matmul(np.matmul(np.linalg.pinv(np.matmul(X.T, X)), X.T), Y)
print(T)
print(J(T, X, Y))
\end{lstlisting}
